% P03.1 ordinary-locally-split LFS descent proof, live integration version.

\begin{proof}[\Cref{thm:p031-lfs-fiber-maximum}]
Fix $x\in X$ and a splitting neighborhood $U$ as in
\eqref{eq:p031-ordinary-local-splitting}.  Write
$\sigma_i:=\pi_i^{-1}:U\to U_i$.  The local description of
locally Fubini--Study functions in \cite[\S7.9]{GJR26} is invariant under
analytic isomorphisms: composing the analytic units in a local presentation
with $\sigma_i$ again gives analytic units.  Hence
\[
                         u_i:=u|_{U_i}\circ\sigma_i
\]
is locally Fubini--Study on $U$ for every $i$.

For $y\in U$, the points
$\sigma_1(y),\ldots,\sigma_d(y)$ are exactly the fiber $\pi^{-1}(y)$, so
\begin{equation}
\label{eq:p031-local-fiber-maximum}
                    \pi_{\max}u|_U=\max_{1\leq i\leq d}u_i.
\end{equation}
Locally Fubini--Study functions are closed under finite maxima, again by
the full real-coefficient local calculus of \cite[\S7.9]{GJR26}.  Thus
$\pi_{\max}u$ is locally Fubini--Study on $U$.  Such splitting
neighborhoods cover $X$, proving the theorem.
\end{proof}

\begin{proof}[\Cref{thm:p031-lfs-envelope-pullback}]
Pullback preserves the local presentation of a locally Fubini--Study
function.  Consequently every base competitor $b$ gives the cover
competitor $b\circ\pi$, and hence
\begin{equation}
\label{eq:p031-lfs-envelope-easy-inequality}
              s_{\mathrm{LFS},E}
                    \geq s_{\mathrm{LFS},X}\circ\pi.
\end{equation}

Conversely, let $u$ be a cover competitor and set
$b:=\pi_{\max}u$.  By \Cref{thm:p031-lfs-fiber-maximum}, the function $b$
is locally Fubini--Study on $X$.  For $x\in X$, every
$e\in\pi^{-1}(x)$ satisfies $u(e)\leq f(x)$; therefore $b\leq f$ and $b$
is a base competitor.  For every $e\in E$ one then has
\[
              u(e)\leq b(\pi(e))
                    \leq s_{\mathrm{LFS},X}(\pi(e)).
\]
Taking the supremum over all cover competitors gives the reverse of
\eqref{eq:p031-lfs-envelope-easy-inequality} whenever the cover family is
nonempty.

Finally, a base competitor pulls back to a cover competitor, while a cover
competitor produces a base competitor by fiberwise maximization.  Thus the
two families are empty simultaneously, and the asserted equality also holds
under the stated empty-family convention.
\end{proof}
